\section{Отображения}

\begin{definition}{}{}
    Пусть заданы множества $X$ и $Y$. Если каждому элементу $x \in X$ сопоставлен единственный элемент $y \in Y$,
    то говорят, что определено \textbf{отображение} $f \colon X \to Y$, действующее из $X$ в $Y$.
\end{definition}

\begin{description}
    \item [Область определения] $\Def(f)$ или $\D(f)$: Множество $X$ называется областью определения отображения.
    \item [Область значений:] Множество $Y$ называется областью значений отображения.
    \item [Образ] $\Im(f)$: Множество всех элементов $y \in Y$, для каждого из которых существует такой $x \in X$,
            что $y = f(x)$, называется образом отображения, причем $\Im(f) \subset Y$,
            где $f$ --- это правило, по которому осуществляется сопоставление элементов.
\end{description}

Отображения бывают:
\begin{description}
    \item[Инъективное] (в $Y$): разные элементы отображаются в разные ---
        если $x_1 \neq x_2$, то $f(x_1) \neq f(x_2)$.
    \item[Сюръективное] (на $Y$): имеет своим образом все множество $Y$ ---
        для каждого $y \in Y$ существует такой $x \in X$, что $y = f(x)$, то есть $\Im(f) = Y$.
    \item[Биективное] (взаимно однозначное соответствие): является одновременно и инъективным и сюръективным ---
        для каждого $y \in Y$ существует такой единственный $x \in X$, что $y = f(x)$.
\end{description}

\begin{definition}{График}{}
    \textbf{Графиком отображения} $f \colon X \to Y$ называется подмножество $X \times Y$,
    определяемое следующим образом:
    $\Gr(f) = \{ (x, y) \in X \times Y \colon y = f(x) \}$.
\end{definition}
